% Figure: a wheel being pulled over a step (curb). FBD shows the pull P at
% the axle, the weight W, and the single contact reaction at the step edge
% once the ground contact lifts. Numbers match the wheel-over-step worked
% example.
\begin{figure}[htbp]
\centering
\begin{tikzpicture}[
  force/.style={draw=engblue, -{Stealth}, very thick},
  rxn/.style={draw=engred, -{Stealth}, very thick},
  scale=1.0,
]
% ground and step
\draw[draw=engblack, very thick] (-0.2,0) -- (3.0,0);
\draw[draw=engblack, very thick] (3.0,0) -- (3.0,1.1) -- (6.2,1.1);
\foreach \x in {0.1,0.5,...,2.9} {\draw[draw=enggray, thin] (\x,0) -- (\x-0.2,-0.28);}
% wheel: radius r, resting against the step edge E at (3.0,1.1)
% centre chosen so wheel touches step corner
\coordinate (E) at (3.0,1.1);
\coordinate (C) at (2.05,1.6);
\def\rw{1.35}
\draw[draw=engblack, very thick] (C) circle (\rw);
\fill[engblack] (C) circle (0.05);
% line from centre to step edge (the contact-reaction line)
\draw[draw=enggray, thin, dashed] (C) -- (E);
% weight down at centre
\draw[force] (C) -- ++(0,-1.7) node[anchor=north, font=\small, engblue] {$W$};
% pull P horizontal at axle
\draw[force] (C) -- ++(1.9,0) node[anchor=west, font=\small, engblue] {$P$};
% contact reaction at step edge, along E->C
\draw[rxn] (E) -- ++($(C)-(E)$) node[anchor=south east, font=\small, engred] {$R$};
% geometry marks
\node[font=\scriptsize, anchor=east] at (2.05,1.6) {$C$};
\node[font=\scriptsize, anchor=north west] at (3.05,1.1) {$E$};
\draw[draw=enggray, thin, <->] (3.0,-0.05) -- (3.0,1.05);
\node[font=\scriptsize, enggray, anchor=west] at (3.08,0.5) {step $h$};
\end{tikzpicture}
\caption[A wheel pulled over a step]{The chosen body is the wheel at the
instant its ground contact lifts, so the only contact left is at the step
corner $E$. A smooth corner can push only along the line from the corner to
the wheel centre $C$, which fixes the direction of the step reaction $R$
before any arithmetic. Three forces act: the weight $W$ down at the axle, the
horizontal pull $P$ at the axle, and the corner reaction $R$ along $EC$. The
three lines meet at $C$, so the system is concurrent and a moment taken about
$E$ kills $R$ at a stroke, leaving one equation for the pull that just starts
the wheel over the step. A taller step needs a larger pull, and a pull
applied at the top of the wheel rather than the axle needs less, which is why
a hand cart is tilted back to climb a curb.}
\label{fig:vol03ch01:wheel-step}
\end{figure}
